\subsection{Explosion and low-variation diagnostics}
\label{sec:generator-diagnostics}

Unrestricted feedback can amplify observations over repeated execution.
In the recorded four-family experiment, the original sampler crosses the
operational explosion criterion in 2,658 of 5,000 trajectories by step 384
(53.16\%). A pretraining prior dominated by such behaviour would allocate much
of its capacity to rapidly growing recurrences. For the intended benchmark,
explosive cases can instead form an explicitly controlled difficulty group,
alongside mechanisms producing sustained variation at usable scales.

\paragraph{Common evaluation procedure.}
The executed notebook
\path{experiments/abstract_syntax_tree/6_stability_methods_tradeoff.ipynb}
uses root seed 20261001, 5,000 independent mechanisms per family, and one
trajectory per mechanism. Each run has one channel, a lag range of 1--12,
and diagnostic horizons $h\in\{96,192,384\}$. The original trees contain
1--6 binary operations, addition probability 0.8, parameter/lag leaf
probabilities 0.55/0.45, and coefficient magnitudes in $[0.05,0.8]$ with
random signs. Explicit time leaves are disabled in this comparison.

The evaluator constructs the history according to each family, generates one
path up to the largest horizon, and reuses its prefixes. It stops at the first
non-finite value or magnitude greater than 100, retaining that failing value.
The event remains counted at every later horizon, including a failure at the
last requested step. Every attempted mechanism stays in the denominator.
For multivariate use, any failing channel flags the system. These rules are
implemented in \path{src/tsfmxai/evaluation/stability.py} and
\path{src/tsfmxai/evaluation/generator_tradeoff.py}.

\paragraph{Low amplitude near zero: the measured diagnostic.}
A completed, non-explosive prefix receives this label when every channel's
last 24 values lies in $[-0.15,0.15]$. The tolerance is three times the
innovation standard deviation used in this experiment, $\sigma_\varepsilon=0.05$.
Thus the label measures occupation of a small band during a finite window.
It can also be produced by a stationary process fluctuating near zero from
the beginning. Convergence of a trajectory to zero is a stronger,
infinite-time property requiring separate evidence.

\paragraph{Nearly constant tails: an additional diagnostic specification.}
To capture stabilization around levels other than zero, let
$\bar x^{\mathrm{tail}}_{k,h}$ be the mean of channel $k$ over the last
$w_{\mathrm{tail}}$ generated observations at horizon $h$. Define
\begin{equation}
 V_h^{\mathrm{tail}}=
 \max_k\max_{h-w_{\mathrm{tail}}<v\leq h}
 \left|x_{k,\tau+v}-\bar x^{\mathrm{tail}}_{k,h}\right|.
 \label{eq:generator-flat-tail}
\end{equation}
A threshold $\delta_{\mathrm{flat}}$ then selects nearly constant tails
around their own mean. Record the mean level, the tolerance, and the innovation
scale together. A corresponding earlier window establishes whether variation
actually decreased during generation; repeated terminal windows provide a
finite-duration persistence check. These additions distinguish an initially
quiet series, flattening towards a nonzero level, and declining amplitude
near zero. This criterion specifies the next extension of the evaluator;
the executed results below report the existing explosion and near-zero-band
diagnostics.

\paragraph{Interpretation and timing.}
An amplitude threshold is an operational safeguard: Gaussian innovations have
unbounded support and can cross any finite threshold even under a stationary
law. The evaluator records finite threshold crossings separately from
non-finite outputs. Timing includes sampling/certification, initialization
(including burn-in), innovation generation, AST execution, and diagnostics.
Completed full-horizon trajectories per second uses the total wall-clock
cost, including failed attempts. This prevents early-stopping generators
from appearing efficient solely because they evaluate fewer time steps.
